\chapter{The Optical System and Myopia Model}
\label{ch:optical_system}

\section{Reduced-eye schematic}

Every model in this report represents the human eye with the standard \emph{reduced
eye} simplification (Gullstrand--Emsley convention): the cornea, aqueous humour,
pupil, and crystalline lens are collapsed into a single equivalent thin lens of power
$P_0$, imaging onto a single retinal plane at posterior distance $v_r$. This is a
deliberate, explicitly documented modeling choice, not an oversight --- it is what
makes a single number $P_0$ meaningful at all, but it means no anterior-chamber
distance, no cornea/lens power split (roughly $2/3$ vs.\ $1/3$ of $P_0$ in a real eye),
and no off-axis retinal curvature are represented. Table~\ref{tab:optical-path} lists
the complete optical path this project's pipeline models.

\begin{table}[htbp]
\centering
\begin{tabular}{llll}
\toprule
Element & Purpose & Refractive index & Role in this model \\
\midrule
Pinhole (opaque film) & Angular filtering & opaque (n/a) & the only design freedom \\
Air gap $d_{gap}$ & Near-field propagation & $n=1.000$ & first ASM propagation stage \\
Cornea + aqueous & Primary refraction & $n\approx1.336$--$1.376$ & absorbed into $P_0$ \\
Pupil (iris aperture) & Variable stop & n/a & clips the effective aperture \\
Crystalline lens & Refracting element & $n\approx1.386$--$1.406$ & absorbed into $P_0$ (fixed) \\
Vitreous humour & Fills posterior chamber & $n'=1.336$ & image-side medium for $v_r$ \\
Retina & Image/sensing plane & n/a & observation plane, $z=v_r$ \\
\bottomrule
\end{tabular}
\caption{Complete optical path from pinhole to retina, and how each element enters the reduced-eye model.}
\label{tab:optical-path}
\end{table}

\section{Myopia as axial elongation}

This project models myopia as purely \emph{axial}: the eye's optical power $P_0$ is
held fixed at its emmetropic (non-myopic) value for every degree of myopia $M$, and
only the eye's physical axial length grows to push the retina past the fixed focal
point. This is the defining assumption of an axial myopia model, as opposed to a
\emph{refractive} one (where $P_0$ itself would change and axial length would stay
fixed); the codebase's own \texttt{myopia\_axial.py} documents this choice explicitly
since no earlier stage in the pipeline had established a convention to match.

Let $n'$ be the vitreous refractive index and $P_0$ the eye's fixed emmetropic power
(diopters). The emmetropic axial length and the myopic axial length at excess power
$M$ (diopters) are
\begin{align}
L_0 &= \frac{n'}{P_0} \label{eq:emmetropic-length}\\
L_{ax}(M) &= \frac{n'}{P_0 - M} \label{eq:myopic-length}
\end{align}
with the eye's fixed focal length $f_{eye}=L_0$ used at \emph{every} myopia level, and
the retina-plane propagation distance for a given $M$ set to $v_r = L_{ax}(M)$. At
default values $P_0=60\,$D, $n'=1.336$: $L_0 \approx 22.27\mm$ ($f_{eye}$), and, e.g.,
at $M=4\,$D, $L_{ax}(4) = 1.336/56 \approx 23.86\mm$ --- the retina sits
$\Delta L \approx 1.59\mm$ further back than the emmetropic focal plane. Equation
\eqref{eq:myopic-length} is singular as $M\to P_0$ and undefined beyond it; the model
raises an explicit error rather than silently returning a negative or infinite length
in that regime, which is never reached at any myopia level swept in this report
($M\le8\,$D $\ll P_0=60\,$D).

\textbf{Refractive-index convention.} Two conventions are supported: with
\texttt{refractive\_index\_correction=True} (this project's default throughout), all
retina-side propagation distances and wavelengths are scaled by the vitreous's own
refractive index $n'=1.336$, the anatomically-scaled convention. With it disabled,
$n'$ is forced to $1.0$ (an air-equivalent convention). \textbf{This one convention
choice, and precisely how it must be threaded consistently through every downstream
calculation, is the exact site of the single most consequential bug found in this
project's history} --- see Chapter~\ref{ch:physics_audit}.

\section{Diopters and the far point}

By definition $D = 1/f$ (focal length in meters). For an eye with power in excess of
emmetropic by $M$ diopters, the far point --- the object distance at which the
\emph{unaccommodated} eye focuses exactly on the retina --- is finite:
\begin{equation}
u_{far} = \frac{1}{M}.
\end{equation}
This project models only unaccommodated, distance-vision viewing (the reduced eye's
$f$ is fixed throughout every chapter), consistent with the ``distant objects'' clause
of the governing problem statement; near-vision/accommodated performance is not
modeled anywhere in this report.
